\setcounter{section}{8}

  \zwischenueberschrift{Teorema pemetaan implisit dan manifold}

Permukaan bola satuan, yang dalam kuliah sebelumnya kita bahas sebagai contoh pemotivasi sebuah manifold, merupakan serat di atas $1$ dari pemetaan diferensiabel
\maabbeledisp {} { \R^3 } { \R
} { (x,y,z) } { x^2+y^2+z^2
} {,}
Pemetaan ini
\definitionsverweis {reguler}{}{}
kecuali di titik nol. Dalam situasi ini, teorema pemetaan implisit memberikan pernyataan yang luas tentang bentuk lokal serat suatu pemetaan
\maabb {\varphi} {\R^n} { \R^m
} {,}
yakni bahwa secara lokal terdapat homeomorfisme antara serat di suatu titik reguler dan sebuah himpunan terbuka di $\R^k$, dengan $k$ selisih antara dimensi ruang asal dan dimensi ruang sasaran. Sebentar lagi akan kita lihat bahwa serat-serat semacam itu bukan hanya manifold topologis, melainkan juga manifold diferensiabel. Kita merumuskan teorema pemetaan implisit dalam suatu versi yang memperlihatkan bahwa serat-serat reguler merupakan manifold diferensiabel.

__NOEDITSECTION__

\inputfaktbeweis
{Satz über implizite Abbildungen/Globale Diffeomorphismen/Induzierte Diffeomorphismen zwischen Faser und Achsenräumen/Fakt}
{Teorema}
{}
{

\faktsituation {Misalkan

\mavergleichskette
{\vergleichskette
{ G
}
{ \subseteq }{ \R^n
}
{  }{
}
{    }{
}
{   }{
}
}
{}{}{}
\definitionsverweis {terbuka}{}{}
dan misalkan
\maabbdisp {\varphi} { G } {\R^m
} {}
suatu
\definitionsverweis {pemetaan terdiferensialkan secara kontinu}{}{.}
Misalkan

\mavergleichskette
{\vergleichskette
{ Z
}
{ = }{ \varphi^{-1}(0)
}
{  }{
}
{    }{
}
{   }{
}
}
{}{}{}
merupakan
\definitionsverweis {serat}{}{}
di atas

\mavergleichskette
{\vergleichskette
{ 0
}
{ \in }{ \R^m
}
{  }{
}
{    }{
}
{   }{
}
}
{}{}{,}
dan misalkan $\varphi$
\definitionsverweis {reguler}{}{}
di setiap titik serat tersebut.}
\faktfolgerung {Maka untuk setiap titik

\mavergleichskette
{\vergleichskette
{ P
}
{ \in }{ Z
}
{  }{
}
{    }{
}
{   }{
}
}
{}{}{}
terdapat suatu lingkungan terbuka

\mavergleichskette
{\vergleichskette
{ P
}
{ \in }{ W
}
{ \subseteq }{ G
}
{    }{
}
{   }{
}
}
{}{}{,}
himpunan-himpunan terbuka

\mavergleichskette
{\vergleichskette
{ V
}
{ \subseteq }{ \R^{n- m}
}
{  }{
}
{    }{
}
{   }{
}
}
{}{}{}
dan

\mavergleichskette
{\vergleichskette
{ V'
}
{ \subseteq }{ \R^{ m }
}
{  }{
}
{    }{
}
{   }{
}
}
{}{}{,}
serta suatu
$C^1$-\definitionsverweis {difeomorfisme}{}{}
\maabbdisp {\theta} { W } { V \times V'
} {}
dengan

\mavergleichskette
{\vergleichskette
{ \varphi\!\mid_{W}
}
{ = }{ p_2 \circ \theta
}
{  }{
}
{    }{
}
{   }{
}
}
{}{}{,}
yang menginduksi bijeksi antara
\mathkor {} {Z \cap W} {dan} {V \times \{0\}} {}
sedemikian sehingga
\definitionsverweis {diferensial total}{}{}

\mathl{\left(D\theta\right)_{Q}}{} untuk setiap

\mavergleichskette
{\vergleichskette
{ Q
}
{ \in }{ W
}
{  }{
}
{    }{
}
{   }{
}
}
{}{}{}
menginduksi bijeksi antara
\mathkor {} {\operatorname{kern}  \left(D\varphi\right)_{ Q }} {dan} {\R^{n-m }} {}
.}
\faktzusatz {}
\faktzusatz {}

}
{

Pernyataan ini telah dibuktikan sekaligus dalam [[Satz über implizite Abbildungen/R/Fakt/Beweis|bukti]]
[[Satz über implizite Abbildungen/R/Fakt|teorema pemetaan implisit]].
Pernyataan tambahan diperoleh dari

\mavergleichskettedisp
{\vergleichskette
{ \operatorname{kern}  \left(D\varphi\right)_{Q}
}
{ \cong} { \operatorname{kern}  \left(Dp_2 \right)_{ \theta(Q)}
}
{ =} { \R^{n- m } \times \{0\} \cong \R^{n-m}
}
{ } {
}
{ } {
}
}
{}{}{.}

Catatan edisi: Sumber mengidentifikasi kernel proyeksi kedua langsung dengan ruang berdimensi lebih rendah. Bentuk di atas menampilkan dahulu subruang yang sebenarnya, lalu isomorfisme kanonisnya.

}

Dari hasil ini, serat itu sendiri memperoleh struktur manifold. Dengan demikian, teorema pemetaan implisit menyediakan bagi kita suatu kelas manifold yang sangat besar. Manifold-manifold ini disebut \stichwort {submanifold tertutup} {,} yang segera akan kita bahas secara lebih sistematis setelah ruang singgung tersedia.

__NOEDITSECTION__

\inputfaktbeweis
{Satz über implizite Abbildungen/Faser ist Mannigfaltigkeit/Fakt}
{Teorema}
{}
{

\faktsituation {Misalkan

\mavergleichskette
{\vergleichskette
{ G
}
{ \subseteq }{ \R^n
}
{  }{
}
{    }{
}
{   }{
}
}
{}{}{}
\definitionsverweis {terbuka}{}{}
dan misalkan
\maabbdisp {\varphi} { G } { \R^{ m }
} {}
suatu
\definitionsverweis {pemetaan terdiferensialkan secara kontinu}{}{.}
Misalkan

\mavergleichskette
{\vergleichskette
{ Z
}
{ = }{ \varphi^{-1}( Q)
}
{  }{
}
{    }{
}
{   }{
}
}
{}{}{}
merupakan
\definitionsverweis {serat}{}{}
di atas suatu titik

\mavergleichskette
{\vergleichskette
{ Q
}
{ \in }{ \R^{ m }
}
{  }{
}
{    }{
}
{   }{
}
}
{}{}{.}}
\faktvoraussetzung {Misalkan
\definitionsverweis {diferensial total}{}{}

\mathl{\left(D\varphi\right)_{ P }}{}
\definitionsverweis {surjektif}{}{}
untuk setiap titik

\mavergleichskette
{\vergleichskette
{ P
}
{ \in }{ Z
}
{  }{
}
{    }{
}
{   }{
}
}
{}{}{.}}
\faktfolgerung {Maka $Z$ merupakan
\definitionsverweis {manifold diferensiabel}{}{}
kelas $C^1$ dengan
\definitionsverweis {dimensi}{}{}

\mathl{n - { m }}{.}}
\faktzusatz {}
\faktzusatz {}

}
{

Dengan mengurangkan nilai serat dari pemetaan, kita dapat mereduksi ke kasus

\mavergleichskette
{\vergleichskette
{ Q
}
{ = }{ 0
}
{  }{
}
{    }{
}
{   }{
}
}
{}{}{}
.
Catatan edisi: Sumber langsung menetapkan nilai serat sama dengan nol; kalimat sebelumnya menambahkan translasi kodomain yang diperlukan agar reduksi itu sah.
Berdasarkan
[[Satz über implizite Abbildungen/Globale Diffeomorphismen/Induzierte Diffeomorphismen zwischen Faser und Achsenräumen/Fakt|Teorema 8.1]],
untuk setiap titik

\mavergleichskette
{\vergleichskette
{ P
}
{ \in }{ Z
}
{  }{
}
{    }{
}
{   }{
}
}
{}{}{}
terdapat suatu lingkungan terbuka

\mavergleichskette
{\vergleichskette
{ P
}
{ \in }{ W
}
{ \subseteq }{ G
}
{    }{
}
{   }{
}
}
{}{}{}
dan suatu
$C^1$-\definitionsverweis {difeomorfisme}{}{}
\maabbdisp {\theta} { W } { V \times V'
} {}
dengan himpunan-himpunan terbuka
\mathkor {} {V \subseteq \R^{n - m }} {dan} {V' \subseteq \R^m} {}
sedemikian sehingga $\theta$ menginduksi bijeksi antara
\mathl{Z \cap W}{} dan

\mavergleichskette
{\vergleichskette
{ V \times \{0\}
}
{ \cong }{ V
}
{  }{
}
{    }{
}
{   }{
}
}
{}{}{}
.
Pembatasan difeomorfisme ini masing-masing pada
\mathl{Z \cap W}{} dan $V$ kita ambil sebagai peta untuk $Z$.
\teilbeweis {}{}{}
{Untuk menunjukkan bahwa hal ini mendefinisikan struktur diferensiabel pada $Z$, misalkan diberikan lingkungan-lingkungan terbuka
\zusatzklammer {di $\R^n$} {} {}
\mathkor {} {W_1} {dan} {W_2} {}
dari

\mavergleichskette
{\vergleichskette
{ P
}
{ \in }{ Z
}
{  }{
}
{    }{
}
{   }{
}
}
{}{}{}
bersama difeomorfisme
\maabbdisp {\theta_1} {W_1 } { V_1 \times V_1'
} {} dan \maabbdisp {\theta_2} {W_2 } { V_2 \times V_2'
} {.}
Dengan beralih ke

\mavergleichskette
{\vergleichskette
{ W
}
{ = }{ W_1 \cap W_2
}
{  }{
}
{    }{
}
{   }{
}
}
{}{}{}
kita dapat mengasumsikan bahwa kedua himpunan terbuka tersebut sama. Pemetaan transisi
\mathl{\theta_2 \circ \theta_1^{-1}}{} merupakan difeomorfisme $C^1$ antara
\zusatzklammer {himpunan-himpunan bagian terbuka dari} {} {}
\mathkor {} {\theta_1(W)} {dan} {\theta_2(W)} {,}
yang memetakan
\mathl{\theta_1(W\cap Z)}{} ke
\mathl{\theta_2(W\cap Z)}{}. Oleh karena itu, menurut
[[Diffeomorphismus/Euklidische Produktmengen/Eingeschränkter Diffeomorphismus/Aufgabe|Soal 52 (Analisis (Osnabrück 2021--2023))]],
pembatasan pemetaan transisi pada himpunan-himpunan bagian tersebut juga merupakan difeomorfisme $C^1$
\zusatzklammer {antara himpunan-himpunan bagian terbuka di \mathlk{\R^{n - m }}{}} {} {.}}
{}

Catatan edisi: Sumber menyebut seluruh irisan sumbu pada kedua kodomain. Setelah lingkungan dipersempit menjadi irisan kedua lingkungan awal, domain yang tepat ialah kedua citra terbuka dari irisan tersebut sebagaimana ditulis di atas.

}

  \zwischenueberschrift{Pemetaan diferensiabel}

\inputdefinition
{}
{

Misalkan
\mathkor {} {L} {dan} {M} {}
dua
\definitionsverweis {manifold}{}{}
$C^k$ dengan
\definitionsverweis {atlas}{}{}
\mathkor {} {(U_i,U_i', \alpha_i, i \in I)} {dan} {(V_j,V_j', \beta_j, j \in J)} {.}
Misalkan

\mavergleichskette
{\vergleichskette
{ 1
}
{ \leq }{ \ell
}
{ \leq }{ k
}
{    }{
}
{   }{
}
}
{}{}{.}
Suatu
\definitionsverweis {pemetaan kontinu}{}{}
\maabbdisp {\varphi} { L } { M
} {}
disebut suatu $C^\ell$-\definitionswort {pemetaan diferensiabel}{,} jika untuk semua
\mathkor {} {i \in I} {dan semua} {j \in J} {}
pemetaan
\maabbdisp {\beta_j \circ  \varphi \circ (\alpha_i)^{-1}} {\alpha_i( \varphi^{-1}(V_j) \cap U_i) } { V_j'
} {}
bersifat
$C^\ell$-\definitionsverweis {diferensiabel}{}{.}

}

Karena
\mathl{\alpha_i \left(  \varphi^{-1}(V_j) \cap U_i  \right)}{} terbuka, definisi ini mereduksi konsep diferensiabilitas bagi pemetaan antara manifold pada konsep diferensiabilitas bagi pemetaan antara himpunan-himpunan terbuka di ruang vektor real. Karena manifold $C^k$ dapat dipandang sebagai manifold $C^\ell$ untuk

\mavergleichskette
{\vergleichskette
{ \ell
}
{ \leq }{ k
}
{  }{
}
{    }{
}
{   }{
}
}
{}{}{}
,
pada dasarnya cukup membicarakan pemetaan $C^k$ antara manifold-manifold $C^k$. Kasus yang khususnya penting adalah
\mathl{k=1,2, \infty}{.} Perhatikan bahwa untuk

\mavergleichskette
{\vergleichskette
{ k
}
{ = }{ 1
}
{  }{
}
{    }{
}
{   }{
}
}
{}{}{}
kita berbicara tentang pemetaan diferensiabel, meskipun
\zusatzklammer {sejauh ini} {} {}
belum ada suatu \anfuehrung{turunan}{.}

__NOEDITSECTION__

\inputfaktbeweis
{Differenzierbare Mannigfaltigkeit/Differenzierbare Abbildungen/Elementare funktorielle Eigenschaften/Fakt}
{Proposisi}
{}
{

\faktsituation {Misalkan
\mathkor {} {L,M} {dan} {N} {}
merupakan
\definitionsverweis {manifold}{}{}
$C^k$.}
\faktuebergang {Maka berlaku pernyataan-pernyataan berikut.}
\faktfolgerung {\aufzaehlungvier{\definitionsverweis {Identitas}{}{}
\maabbdisp {\operatorname{Id} \,} { L } { L
} {}
merupakan
\definitionsverweis {pemetaan}{}{}
$C^k$.
}{Setiap
\definitionsverweis {pemetaan konstan}{}{}
\maabbdisp {\varphi} { L } { M
} {}
merupakan
\definitionsverweis {pemetaan}{}{}
$C^k$.
}{Untuk setiap
\definitionsverweis {himpunan bagian terbuka}{}{}

\mavergleichskette
{\vergleichskette
{ U
}
{ \subseteq }{ L
}
{  }{
}
{    }{
}
{   }{
}
}
{}{}{}
\definitionsverweis {penyematan terbuka}{}{}
\maabb {} { U } { L
} {}
merupakan
\definitionsverweis {pemetaan}{}{}
$C^k$.
}{Misalkan
\maabbdisp {\varphi} { L } { M
} {}
dan
\maabbdisp {\psi} { M } { N
} {}
merupakan
\definitionsverweis {pemetaan}{}{}
$C^k$. Maka
\definitionsverweis {komposisi}{}{}
\maabbdisp {\psi \circ \varphi} { L } { N
} {}
juga merupakan
\definitionsverweis {pemetaan}{}{}
$C^k$.
}}
\faktzusatz {}
\faktzusatz {}

}
{

\teilbeweis {}{}{}
{(1). Pemetaan-pemetaan yang harus diperiksa tepatnya adalah pemetaan transisi
\mathl{\alpha_j \circ \alpha_i^{-1}}{,} yang menurut definisi suatu
\definitionsverweis {manifold diferensiabel}{}{}
$C^k$ merupakan
\definitionsverweis {difeomorfisme}{}{}
$C^k$.}
{}
\teilbeweis {}{}{}
{(2). Dalam setiap peta, pemetaan yang harus diperiksa bersifat konstan, sehingga diferensiabel hingga sebarang orde.}
{}
\teilbeweis {}{}{}
{(3). Pemetaan-pemetaan yang harus diperiksa untuk

\mavergleichskette
{\vergleichskette
{ i,j
}
{ \in }{ I
}
{  }{
}
{    }{
}
{   }{
}
}
{}{}{}
sama dengan

\mathdisp {\alpha_i \left(  U \cap U_i  \cap U_j  \right)  \subseteq  \alpha_i \left(  U_i  \cap U_j  \right) \stackrel{ \alpha_j \circ \alpha_i^{-1} }{  \longrightarrow} U'_j} { , }

sehingga merupakan penyematan terbuka yang diikuti oleh pemetaan transisi diferensiabel.}
{}
\teilbeweis {}{}{}
{(4). Misalkan
\maabbdisp {\gamma_{  \ell }} {W_{  \ell }} {W'_{  \ell }
} {}
merupakan peta-peta untuk $N$. Maka, untuk semua kombinasi indeks yang mungkin, komposisi
\zusatzklammer {yang dibatasi pada himpunan-himpunan bagian terbuka tertentu} {} {}

\mavergleichskettedisphandlinks
{\vergleichskettedisphandlinks
{ \gamma_{  \ell } \circ (\psi  \circ \varphi) \circ \alpha_i^{-1}
}
{ =} { \gamma_{  \ell } \circ  \psi  \circ \beta_j^{-1} \circ \beta_j \circ \varphi \circ \alpha_i^{-1}
}
{ =} { \left(  \gamma_{  \ell } \circ \psi \circ \beta_j^{-1}  \right)  \circ \left(  \beta_j \circ \varphi \circ \alpha_i^{-1}  \right)
}
{ } {
}
{ } {
}
}
{}{}{}
bersifat
\definitionsverweis {diferensiabel}{}{}
menurut [[Totale Differenzierbarkeit/R/Kettenregel/Fakt|aturan rantai]].
Untuk

\mavergleichskette
{\vergleichskette
{ k
}
{ \geq }{ 2
}
{  }{
}
{    }{
}
{   }{
}
}
{}{}{}
kita menggunakan
[[Stetig differenzierbar/K/Höherdimensional/Mehrfach/Hintereinanderschaltung/Aufgabe|Soal 46.10 (Analisis (Osnabrück 2021--2023))]].}
{}

}

\inputdefinition
{}
{

Misalkan
\mathkor {} {L} {dan} {M} {}
dua
\definitionsverweis {manifold}{}{}
$C^k$. Suatu
\definitionsverweis {homeomorfisme}{}{}
\maabbdisp {\varphi} { L } { M
} {}
disebut
\definitionswortpraemath {C^k}{ difeomorfisme }{,}
jika
\mathkor {} {\varphi} {dan} {\varphi^{-1}} {}
keduanya merupakan
\definitionsverweis {pemetaan}{}{}
$C^k$.

}

\inputdefinition
{}
{

Dua
\definitionsverweis {manifold}{}{}
$C^k$
\mathkor {} {L} {dan} {M} {}
disebut
\definitionswortpraemath {C^k}{ difeomorfik }{,}
jika terdapat suatu
\definitionsverweis {difeomorfisme}{}{}
$C^k$ di antara keduanya.

}

\inputbemerkung
{}
{

Untuk suatu
\definitionsverweis {manifold}{}{}
$C^k$ $M$ dengan atlas $C^k$
\mathl{\left(  U_i,U_i',\alpha_i, i \in I  \right)}{} terdapat suatu \stichwort {atlas maksimal} {,} yang kompatibel dengan struktur diferensiabel yang diberikan oleh atlas tersebut. Atlas ini terdiri atas himpunan semua
\definitionsverweis {homeomorfisme}{}{}
\maabbdisp {\beta} { U } { V
} {}
dengan himpunan-himpunan terbuka

\mavergleichskette
{\vergleichskette
{ U
}
{ \subseteq }{ M
}
{  }{
}
{    }{
}
{   }{
}
}
{}{}{}
dan

\mavergleichskette
{\vergleichskette
{ V
}
{ \subseteq }{ \R^n
}
{  }{
}
{    }{
}
{   }{
}
}
{}{}{}
yang mempunyai sifat bahwa pemetaan ini beserta inversnya merupakan
\definitionsverweis {pemetaan}{}{}
$C^k$
\zusatzklammer {terhadap struktur yang diberikan oleh atlas tersebut} {} {.}
Atlas maksimal ini tentu memuat atlas awal, tetapi secara umum jauh lebih besar. Sebagai contoh, untuk setiap peta
\maabb {\beta} { U } { V
} {}
dan setiap himpunan bagian terbuka

\mavergleichskette
{\vergleichskette
{ U'
}
{ \subseteq }{ U
}
{  }{
}
{    }{
}
{   }{
}
}
{}{}{}
atlas tersebut juga memuat pembatasan peta itu pada $U'$. Hal yang penting adalah bahwa pemetaan identitas
\maabbdisp {\operatorname{Id}} {(M,A)} {(M,B)
} {,}
dengan $A$ menyatakan atlas awal dan $B$ atlas maksimal, merupakan
\definitionsverweis {difeomorfisme}{}{}
$C^k$ antara kedua manifold tersebut, sebagaimana langsung mengikuti dari definisi. Struktur diferensiabel pada manifold yang direpresentasikan oleh atlas lebih penting daripada atlas itu sendiri; struktur ini menentukan pemetaan mana yang diferensiabel dan pemetaan mana yang merupakan difeomorfisme. Peta-peta dalam atlas maksimal kadang-kadang juga disebut peta
\zusatzklammer {yang diperumum} {} {}
bagi manifold tersebut.

Catatan edisi: Sumber hanya mensyaratkan regularitas pemetaan maju. Regularitas invers ditambahkan karena sebuah homeomorfisme yang diferensiabel belum tentu merupakan difeomorfisme dan belum tentu kompatibel dengan atlas awal.

}

  \zwischenueberschrift{Fungsi diferensiabel}

Suatu
\definitionsverweis {pemetaan diferensiabel}{}{}
$C^k$
\maabbdisp {f} {M} {\R
} {}
dari suatu
\definitionsverweis {manifold diferensiabel}{}{}
ke bilangan real juga disebut \stichwort {fungsi diferensiabel} {} $C^k$. Menurut definisi, ini berarti bahwa untuk setiap peta
\maabbdisp {\alpha} {U} {V
} {}
fungsi komposit
\maabbdisp {f \circ \alpha^{-1}} {V} {\R
} {}
merupakan
\definitionsverweis {fungsi}{}{}
$C^k$. Himpunan semua fungsi $C^k$ pada $M$ dinotasikan dengan
\mathl{C^k(M,\R)}{}.

__NOEDITSECTION__
\inputfaktbeweis
{Mannigfaltigkeit/Differenzierbare Funktionen/Strukturelle Eigenschaften/Fakt}
{Lema}
{}
{

\faktsituation {Misalkan $M$ suatu
\definitionsverweis {manifold diferensiabel}{}{} dan
\maabbdisp {f,g} { M } {\R
} {}
merupakan
\definitionsverweis {fungsi diferensiabel}{}{}
pada $M$.}
\faktuebergang {Maka berlaku pernyataan-pernyataan berikut.}
\faktfolgerung {\aufzaehlungvier{Pemetaan
\maabbeledisp {f \times g} { M } {\R^2
} { x } {(f(x),g(x))
} {,} bersifat diferensiabel.
}{ $f+g$
\definitionsverweis {diferensiabel}{}{.}
}{ $f \cdot g$
\definitionsverweis {diferensiabel}{}{.}
}{Jika $f$ tidak mempunyai titik nol, maka
\mathl{f^{-1}}{} juga diferensiabel.
}}
\faktzusatz {}
\faktzusatz {}

 }
{ Lihat [[Mannigfaltigkeit/Differenzierbare Funktionen/Strukturelle Eigenschaften/Fakt/Beweis/Aufgabe|Soal 8.3]]. }

Secara khusus, fungsi-fungsi diferensiabel pada suatu manifold membentuk
\definitionsverweis {gelanggang komutatif}{}{.}

Jika
\maabbdisp {\alpha} {U} {V
} {}
merupakan suatu peta dengan

\mavergleichskette
{\vergleichskette
{ V
}
{ \subseteq }{ \R^n
}
{  }{
}
{    }{
}
{   }{
}
}
{}{}{}
terbuka, maka setiap proyeksi $x_i$ menghasilkan suatu fungsi diferensiabel
\maabbdisp {x_i \circ \alpha} {U} {\R
} {,}
yang biasanya kembali dinotasikan dengan $x_i$. Dalam hal ini dikatakan bahwa fungsi-fungsi
\mathl{x_1 , \ldots , x_n}{} membentuk \stichwort {koordinat diferensiabel} {} untuk

\mavergleichskette
{\vergleichskette
{ U
}
{ \subseteq }{ M
}
{  }{
}
{    }{
}
{   }{
}
}
{}{}{}
.
Bagi suatu fungsi yang terdiferensialkan secara kontinu
\maabbdisp {f} {U} {\R
} {}
menurut definisi fungsi
\maabbdisp {f \circ \alpha^{-1}} {V} {\R
} {}
\definitionsverweis {terdiferensialkan secara kontinu}{}{,}
artinya, untuk setiap $i$ terdapat
\definitionsverweis {turunan parsial}{}{}

\mathdisp {{ \frac{   \partial  (f \circ \alpha^{-1})  }{  \partial   x_i  } }} { , }

yang kembali merupakan fungsi
\zusatzklammer {kontinu} {} {}
pada $V$. Oleh karena itu,

\mathdisp {{ \frac{   \partial  (f \circ \alpha^{-1})   }{  \partial   x_i  } }  \circ \alpha} {  }

merupakan fungsi pada $U$. Fungsi ini secara umum kembali dinotasikan dengan
\mathl{{ \frac{   \partial  f  }{  \partial   x_i  } }}{}.

 [[Kategorie:Latexseite]]
